% !TeX root = ../组合数学笔记.tex
% 由主文档通过 \input 载入；请编译主文档。
\setcounter{chapter}{7} % 保留原书章号。
\chapter{指数生成函数与集合分拆}
\label{ch:egf}

不同标签会产生额外的分配因子。例如把 $n$ 个人分给两个有名称的队伍，大小分别为 $k,n-k$，必须先选择哪些人进入第一队。指数生成函数把这样的二项式因子吸收到分母中，使标号结构的组合仍能使用乘法。

\section{标签分配与指数型乘法}

\begin{definition}
若 $a_n$ 计数固定标签集 $[n]$ 上的一类结构，其\emph{指数生成函数（exponential generating function）}为 $A(x)=\sum_{n\ge0}a_nx^n/n!$。相同大小的标签集通过重标号携带相同类型结构，且重标号与复合相容。
\end{definition}

设一个结构由两个有序组件组成，标签集互不相交且并为 $[n]$。先选第一组件的 $k$ 个标签，再分别赋予结构，得
\begin{equation}
 c_n=\sum_{k=0}^n\binom nk a_kb_{n-k},\qquad
 \sum_{n\ge0}c_n\frac{x^n}{n!}=A(x)B(x).
 \label{eq:egf-product}
\end{equation}
等式也可由乘积展开逐项验证，因为 $1/(k!(n-k)!)=\binom nk/n!$。注意 $[x^n]A(x)=a_n/n!$，实际对象数还须乘以 $n!$。普通生成函数和指数生成函数可编码同一数列，但其乘法的组合含义不同。

\section{第二类 Stirling 数与 Bell 数}

\begin{definition}
$[n]$ 的\emph{集合分拆（set partition）}是若干非空、两两不交且并为 $[n]$ 的子集所组成的集合，各子集称为块。恰有 $k$ 个块的数量称为第二类 Stirling 数，记为 $S(n,k)$；规定 $S(0,0)=1$，$n>0$ 时 $S(n,0)=0$，其余不可能的参数取零。
\end{definition}

对 $n\ge1$，元素 $n$ 所在块若是单元素块，删除后剩 $S(n-1,k-1)$ 种；否则从该块删除 $n$，得到 $k$ 块分拆，再从 $k$ 个块中选择插入位置。故
\[
 S(n,k)=S(n-1,k-1)+kS(n-1,k).
\]
将 $k$ 个无名称块分别命名为 $1,\ldots,k$，得到满射的纤维，因此满射数为 $k!S(n,k)$。结合容斥，得到 $S(n,k)=\frac1{k!}\sum_{j=0}^k(-1)^j\binom kj(k-j)^n$（$k\ge1$）。

\emph{Bell 数（Bell number）} $B_n=\sum_{k=0}^nS(n,k)$ 计数全部集合分拆，$B_0=1$。在 $[n+1]$ 的分拆中，选择与 $n+1$ 同块的 $n-k$ 个元素，余下 $k$ 个任意分拆，得到 $B_{n+1}=\sum_{k=0}^n\binom nkB_k$。

\begin{example}
四个人分成两组且组无名称，有 $S(4,2)=7$ 种：大小为 $1+3$ 的有 $4$ 种，大小为 $2+2$ 的有 $\binom42/2=3$ 种。分给两个有名称且均非空的实验室，则有 $2!\cdot7=14$ 种。四人的全部分组数为 $B_4=1+7+6+1=15$。
\end{example}

\section{第一类 Stirling 数与置换循环}

一个置换可以唯一分解为不交循环，循环内部的书写起点与循环之间的书写顺序不属于额外结构。令 $c(n,k)$ 计数 $[n]$ 上恰有 $k$ 个循环的置换，称为无符号第一类 Stirling 数；取 $c(0,0)=1$，其他不可能参数为零。

元素 $n$ 或者单独形成一个循环，或者插入某个已有元素之后。前者贡献 $c(n-1,k-1)$，后者有 $n-1$ 个插入位置，删除 $n$ 能唯一恢复原置换。因此对 $n\ge1$，
\[
 c(n,k)=c(n-1,k-1)+(n-1)c(n-1,k),\qquad
 \sum_{k=0}^n c(n,k)t^k=t(t+1)\cdots(t+n-1).
\]
多项式恒等式由递推归纳，$n=0$ 时右侧为空积 $1$。有符号约定另记为 $s(n,k)=(-1)^{n-k}c(n,k)$，不能把它与这里的非负计数混用。例如 $c(n,1)=(n-1)!$（$n\ge1$），因为一个循环可固定最小元素为书写起点，其余任意排列。

\section{指数公式与非空组件}

\begin{theorem}[标号结构的指数公式]
设组件均非空，其指数生成函数为 $C(x)$，故 $C(0)=0$。若整体结构可唯一分解为一组无序、标签互不相交的组件，则恰有 $k$ 个组件的生成函数为 $C(x)^k/k!$，全部结构的生成函数为 $\exp(C(x))$。
\end{theorem}
\begin{proof}
由式\eqref{eq:egf-product}反复相乘，$C(x)^k$ 计数有顺序的 $k$ 个组件。组件的标签集非空且互不相交，故它们作为带标签对象彼此不同；忘掉顺序后，每个无序组件集恰有 $k!$ 个原像，可以等重相除。对 $k$ 求和得到指数函数。固定总大小 $n$ 时至多有 $n$ 个非空组件，因而逐系数求和有限。
\end{proof}

一个非空标签集本身恰有一种块结构，所以 $C(x)=e^x-1$，得到
\begin{equation}
 \sum_{n\ge0}S(n,k)\frac{x^n}{n!}=\frac{(e^x-1)^k}{k!},\qquad
 \sum_{n\ge0}B_n\frac{x^n}{n!}=\exp(e^x-1).
 \label{eq:stirling-bell-egf}
\end{equation}
一个大小为 $j\ge1$ 的置换循环有 $(j-1)!$ 种，因此组件生成函数为 $\sum_{j\ge1}x^j/j=-\log(1-x)$。全部置换的生成函数为 $\exp(-\log(1-x))=1/(1-x)$，系数乘 $n!$ 后即 $n!$。错排禁止一元循环，所以生成函数为 $e^{-x}/(1-x)$，与式\eqref{eq:derangements}一致。这里指数与对数均用其形式级数定义，相应恒等式可由形式微分和常数项唯一性验证。

\begin{example}[单人和配对]
把 $n$ 个不同的人分成单人组或双人组，组件生成函数为 $x+x^2/2$，整体为 $\exp(x+x^2/2)$。若双人组数为 $j$，先把 $2j$ 人配对，再让其他人各自单独成组，数量为 $n!/((n-2j)!2^jj!)$。对 $0\le j\le\lfloor n/2\rfloor$ 求和；$n=4$ 时为 $1+6+3=10$。同样的数也计数只含一元、二元循环的置换。
\end{example}

\section{习题}
\begin{enumerate}
  \item\label{ex:egf-1} 用递推计算 $S(5,2)$、$S(5,3)$，再由满射计数核对前者。
  \item\label{ex:egf-2} 求 $B_5$，并解释递推中 $\binom nk$ 的标签选择含义。
  \item\label{ex:egf-3} 写出所有块大小均为偶数的集合分拆的指数生成函数，求 $n=4$ 时的数量。
  \item\label{ex:egf-4} 写出没有二元循环的置换的指数生成函数；保留一元循环，与错排条件作比较。
\end{enumerate}

